% 问题三·子问题三：M2/P 配比条件面板与 Q2--Q3 接口敏感性
% 本文件为正文片段，遵循 paper/main.tex 的唯一入口约定；不定义文档环境。
% 证据运行：q3-p-conditional-20260925-r01、q2-q3-interface-sensitivity-20260926-r01。

\subsection{问题三子问题三：M2/P 配比条件面板与 Q2--Q3 接口敏感性}
\label{sec:q3-subproblem3}

\subsubsection{配比进入方式}

对 A4/A5 中实际观测的六个配比候选，记配比向量为
$p=(p_1,\ldots,p_J)$，满足
\begin{equation}
 p_j\ge0,\qquad \sum_{j=1}^{J}p_j=1.
\end{equation}
给定 A 侧各域质量值 $q=(q_1,\ldots,q_J)$，先计算
\begin{equation}
 Q_A(p)=q^\top p=\sum_{j=1}^{J}q_jp_j.
 \label{eq:q3-sp3-qa}
\end{equation}
$Q_A(p)$ 只作为 A 侧的条件标签，不直接加入 B 侧 Loss，也不替换
M1 中的原生 $Q_B$。因此，六个候选与 Q3 的 15 个预算--上下文情景组成
90 行条件面板，表示条件组合而非新增观测样本。

\subsubsection{假设质量接口}

A 侧 $Q_A$ 与 B6/B7 原生 $Q_B$ 没有共同量尺和样本级连接键。为进行可计算的
接口压力测试，本文使用明确标注为假设的区间映射
\begin{equation}
 Q_B=0.1+0.9\,\operatorname{clip}
 \left(\frac{Q_A-q_{\min}}{q_{\max}-q_{\min}},0,1\right),
 \label{eq:q3-sp3-hyp-map}
\end{equation}
其中
\begin{equation}
 q_{\min}=51.7014075324,\qquad q_{\max}=64.3357099536.
 \label{eq:q3-sp3-map-range}
\end{equation}
该式仅将 Q1 候选的 $0$--$100$ 数值压缩到 B6 的 $[0.1,1.0]$ 计算区间，
没有使用 A--B 配对样本估计标度关系。正文对该接口保留三点限制：
\begin{enumerate}
  \item 这不是 $Q_A,Q_B$ 的实证关系；
  \item 在没有共同连接键和标度验证时，不能识别 $G_{\mathrm{bridge}}$；
  \item 映射产生的是反事实条件情景，不是新增观测样本。
\end{enumerate}

\subsubsection{映射情景与计算设计}

并列保留五种 M2 构造。五种值先保留在 Q1 的 $0$--$100$ 尺度，只有进入 Q3
接口压力测试时才使用式~\eqref{eq:q3-sp3-hyp-map}：
\begin{itemize}
  \item 部分映射原值 \texttt{M2\_partial\_raw}；
  \item 部分映射重归一化 \texttt{M2\_partial\_renorm}；
  \item 未映射质量取低端点 \texttt{M2\_interval\_low}；
  \item 未映射质量取中、高端点 \texttt{M2\_interval\_mid} 和
  \texttt{M2\_interval\_high}。
\end{itemize}

接口敏感性矩阵包含
\begin{equation}
 6\times5\times3\times5\times3=1350
\end{equation}
行，分别对应六个候选、五种映射、三类质量成本、五档上下文和三档预算。每行
报告 $Q_A$、假设 $Q_B$、$N^*$、$D^*$、$L^*$、质量成本占比、预算活跃状态、
边界状态和支持范围状态。M0 与 B6 原生 M1 只作独立参考，不参与参数平均，
也不与接口行合并拟合。

\subsubsection{配比条件面板结果}

六个候选的行号为 46、97、170、186、270 和 323。观测标准化 Loss 最低的候选
为 170，而跨模型平均预测排名第一的是 46；这种排序差异说明离散候选面板不能
推出稳定的唯一最优配方。因而本节把 $p$ 作为条件标签与 Q3 的 $N$--$D$ 资源
解并列呈现，不把 $p$ 写入 B 侧目标函数。

\subsubsection{Q2--Q3 接口敏感性结果}

在指数型质量成本、$L_{\mathrm{ctx}}=2048$、$C=10^{22}$ FLOPs 的代表性
情景中，六个候选按映射情景取中位数，结果如表~\ref{tab:q3-sp3-representative}。

\begin{table}[htbp]
  \centering
  \caption{代表性预算--上下文下的 Q2--Q3 接口情景中位数}
  \label{tab:q3-sp3-representative}
  \small
  \begin{tabular}{lrrrrrr}
    \toprule
    映射情景 & $Q_A$ & 假设 $Q_B$ & $N^*$ & $D^*$ & $L^*$ & $\rho_Q$\\
    \midrule
    partial raw & 33.39 & 0.100 & 8.094 & 192.756 & 2.414 & 0.000\\
    interval low & 57.21 & 0.493 & 8.132 & 191.168 & 2.272 & 0.004\\
    interval mid & 61.83 & 0.822 & 8.359 & 181.994 & 2.155 & 0.025\\
    partial renorm & 62.24 & 0.851 & 8.407 & 180.129 & 2.145 & 0.029\\
    interval high & 63.54 & 0.943 & 8.630 & 171.916 & 2.113 & 0.049\\
    \bottomrule
  \end{tabular}
\end{table}

在该代表性场景中，假设 $Q_B$ 增大时，$N^*$ 增大、$D^*$ 减小、质量成本
占比上升而 Loss 下降。全部 1350 行在冻结 M1 和给定成本函数下保持这一条件
方向，但该方向不证明提高 A 侧配比质量必然提高 B6 的 $Q_B$。六个
\texttt{partial\_raw} 候选均被压到 $Q_B=0.1$ 下界，说明把未映射质量直接视为
零贡献会把接口解推向低质量端；这不能解释为 B6 已观测到的低质量事实。

作为独立参照，同类预算--上下文下原生 B6 M1 优化给出
$Q_B=1.0$、$N^*=8.834$、$D^*=164.915$、$L^*=2.094$；该结果不能与
M2/P 假设映射情景合并解释。

\subsubsection{数值复核、稳健性与资源配置变化}

六个候选的 $Q_A$、假设 $Q_B$ 及全部 1350 行输出均为有限值；所有假设
$Q_B$ 位于 B6 的 $[0.1,1.0]$ 支持域，$N,D$ 位于 B1 支持域；最大相对预算
误差为 $9.3193240576\times10^{-12}$。输出表
\texttt{m2\_q3\_fixed\_q\_scenarios.csv} 和
\texttt{m2\_q3\_interface\_summary.csv} 的 SHA256 分别为：
\begin{center}
{\scriptsize\texttt{a5009de0b6cc88d43b61a65dccaec79600bbbbc060f4b610b269d1be9dbb95d4}}\\
{\scriptsize\texttt{f410270c93b2fefed1af2f56463786835740354d701d802e3b503750a150407e}}。
\end{center}
复核同时保持以下接口状态：
\begin{center}
\texttt{a\_b\_row\_join=false},\quad
\texttt{formal\_m3\_fit=false},\quad
\texttt{b8\_used=false}.
\end{center}

因此，接口选择本身会改变资源配置：在代表性情景下，从 partial raw 到
interval high，$N^*$ 由 8.094 增至 8.630，$D^*$ 由 192.756 降至 171.916，
而质量成本占比由 0 增至 0.049。该变化是映射假设造成的条件差异，不是观测到的
配比因果效应。

\subsubsection{结果讨论与不能声称的内容}

在当前证据范围内，可以报告六个已观测配比候选的 $Q_A(p)$ 条件面板、五种
接口构造在不同成本、预算和上下文下对 $Q_B,N^*,D^*,L^*$ 及质量成本占比的
条件影响，以及映射构造造成的资源配置变化。M3-H 有效数据效率模型只能作为
理论反事实接口使用。

不能据此声称 A--B 联合 Loss 已经验证，不能声称 $Q_A$ 与 $Q_B$ 已建立稳定
关系，不能声称 $G_{\mathrm{bridge}}$ 已被识别，不能声称六个候选中存在现实
中的唯一最优配比，也不能把 1350 行情景写成 1350 个新增观测样本或把支持域外
结果写成现实中的全局最优。

\subsubsection{小结}

在 B1、B6/B7 的观测支持域及给定预算、上下文和成本假设下，本小问完成了六个
观测配比候选的条件面板和 $Q_A\to Q_B$ 假设接口敏感性分析。由于 A、B 数据
缺少共同连接键，且 $Q_A$ 与 $Q_B$ 尚未统一标度，本文将配比结果作为条件面板，
将四量模型保留为理论接口，不把条件情景解释为跨来源联合最优。
